\chapter{Metodologia}
\label{sec:avaliacao}

\section{Parâmetro e identificação}

Seja $Y_{it}(1)$ o desfecho do município $i$ no trimestre $t$ sob a lei e
$Y_{it}(0)$ o desfecho sem ela. O parâmetro é o efeito médio sobre os tratados,
\[
ATT(t) = E\big[\,Y_{t}(1) - Y_{t}(0) \,\big|\, G = 2019\,\big], \qquad t \geq 2019,
\]
em que $G = 2019$ identifica os municípios cearenses da Chapada, tratados a
partir do primeiro trimestre de 2019, e a média é ponderada pelo número de
nascimentos. É o efeito de
estar sob a lei, por todos os caminhos da Figura~\ref{fig:dag}, no nível do
município, e não o efeito de deixar de ser pulverizado.

O contrafactual é construído com os municípios potiguares da Chapada, e a
identificação apoia-se em quatro hipóteses: (i) tendências paralelas, isto é,
sem a lei o desfecho dos municípios cearenses teria evoluído como o dos
potiguares; (ii) ausência de antecipação antes do período de base; (iii)
sobreposição, a examinar numa tabela de equilíbrio de linha de base; e (iv) o
grupo de comparação não é afetado pela lei. A Tabela~\ref{tab:ameacas}
confronta as ameaças com o que o desenho faz. A fronteira deixa uma delas
aberta: todos os tratados estão no Ceará, de modo que uma mudança estadual
simultânea se confunde com a lei.

\begin{table}[ht]
\centering\small
\caption{Ameaças à validade interna}
\label{tab:ameacas}
\begin{tabularx}{\textwidth}{@{}p{3.2cm}LL@{}}
\toprule
Ameaça & O que é aqui & Resposta do desenho \\
\midrule
Seca e choques comuns & Seca de 2012--2017; zika; pandemia & Mesma chapada e mesmo clima nos dois lados \\
Choque estadual & Política cearense de saúde próxima de 2019 & Diferença tripla; desenho de calendário \\
Contaminação do controle & Deriva e aquífero cruzam a divisa & Viés contra achar efeito; sensibilidade sem os municípios da divisa \\
Antecipação & Projeto de lei de 2015 & Sensibilidade com o período de base em 2014 \\
Mudança no controle & Drones no lado potiguar & Sensibilidade com o pós-período encurtado \\
Seleção para nascer vivo & Menos perda fetal altera quem nasce & Óbito fetal, razão de sexos e número de nascimentos como desfechos \\
\bottomrule
\end{tabularx}
\fonte{Elaboração própria.}
\end{table}

\section{Estimador}

O estimador é o de \citeonline{callaway2021multiple}, com uma única coorte
tratada (2019) e os municípios potiguares como grupo nunca tratado. Ele entrega
um $ATT(t)$ por período, e os períodos anteriores à lei servem de teste de
tendência. Usa-se a versão duplamente robusta de \citeonline{santanna2020drdid},
com covariáveis de linha de base (população e urbanização em 2010, participação
da agricultura não familiar em 2017) e células ponderadas pelo número de
nascimentos. Com dois municípios tratados no núcleo, e três com Russas, o ajuste por covariáveis tem
pouco suporte, e por isso reporta-se também a versão sem covariáveis. Variáveis
que a lei pode alterar, como pré-natal e tipo de parto, não entram.

A lei entra em vigor em janeiro de 2019, e uma criança nascida em fevereiro teve
a gestação quase toda sob o regime anterior. O efeito deve, portanto, crescer ao
longo de 2019 e estabilizar a partir das concepções daquele ano. A previsão fica
registrada antes da estimação.

\section{Amostra e dados}

Os tratados são os municípios cearenses da Chapada: no núcleo, Limoeiro do Norte
e Quixeré, que concentravam a aplicação por aeronave em 2006; na versão ampla,
também Russas. Os controles são os vinte e um municípios potiguares
da delimitação proposta para a Região Integrada de Desenvolvimento da Chapada do
Apodi. O período principal vai de 2010 a 19 de dezembro de 2024. A unidade é o
município de residência da mãe por trimestre. Não há amostragem: usa-se o
universo dos nascimentos registrados, e as listas de municípios são fixadas
antes de se examinar qualquer desfecho posterior a 2019.

\begin{table}[ht]
\centering\small
\caption{Variáveis e fontes}
\label{tab:dados}
\begin{tabularx}{\textwidth}{@{}p{3.6cm}LL@{}}
\toprule
Papel & Variável & Fonte \\
\midrule
Desfecho primário & Peso ao nascer (g) & SINASC \\
Desfechos secundários & Baixo peso; duração da gestação; prematuridade & SINASC \\
Perda fetal & Óbito fetal; razão de sexos; número de nascimentos & SIM; SINASC \\
Primeiro estágio & Área aplicada por avião, por município e mês & Relatórios ao Ministério da Agricultura (acesso à informação) \\
Substituição & Modalidade de aplicação receitada & Receituário agronômico (acesso à informação) \\
Exposição ocupacional & Intoxicação não intencional por agrotóxico & SINAN; SIH \\
Economia & Área, produção e valor da banana; emprego formal & PAM; RAIS \\
Heterogeneidade & Aeronave por tipo de agricultura; agricultura não familiar & Censo Agropecuário 2006 e 2017 \\
\bottomrule
\end{tabularx}
\fonte{Elaboração própria.}
\end{table}

\section{Robustez e mecanismos}
\label{sec:robustez}

\textbf{Diferença tripla.} Um choque estadual atinge o Ceará inteiro, e a lei só
importa onde havia avião. Compara-se a Chapada com o restante de cada estado, e
a diferença cearense com a potiguar. O restante do Ceará não é de todo isento:
9 dos 36 estabelecimentos com aeronave em 2006 estavam fora de Limoeiro do Norte
e Quixeré, o que atenua o contraste. A identificação exige uma hipótese de
tendências paralelas sobre a diferença entre os dois contrastes
\cite{olden2022triple}.

\textbf{Calendário.} O avião voava na quadra chuvosa, e
\citeonline{calzada2023bananas} encontram o efeito no primeiro e no segundo
trimestre de gestação. Compara-se, dentro de cada município, a coorte de
concepção cuja gestação inicial atravessa a estação de aplicação com as demais,
antes e depois da lei, nos dois lados. A comparação absorve o que varia por
município e ano. O ganho é de identificação; o desenho não aumenta o efeito
esperado.

\textbf{Heterogeneidade.} O Censo de 2006 mostra que o avião era sobretudo de
estabelecimentos não familiares. Testa-se se o efeito é maior onde a agricultura
não familiar pesa mais.

\textbf{Mecanismos.} A substituição de tecnologia mede-se pelo receituário e
pelos registros de operadores; a exposição ocupacional, pelas intoxicações não
intencionais, com a tentativa de suicídio como desfecho placebo; a renda, por
produção e emprego; a composição, pelas características maternas.

\textbf{Cenários do primeiro estágio.} Se os registros mostrarem voo até 2018, o
contraste de 2019 mede o efeito de a lei remover a exposição. Se mostrarem que
o lado cearense parou em 2012, com a seca, o contraste informativo é o de 2012,
sob a hipótese adicional de que a seca afetou os dois lados por igual, e o
pré-período recua a 2005; o de 2019
passa a testar se a lei acrescentou algo. A regra é fixada antes de os registros
serem abertos.

\section{Poder estatístico}
\label{sec:poder}

O efeito de referência é o de \citeonline{calzada2023bananas}: 80 gramas a menos
no peso ao nascer para filhos de mães que moram a até 150 metros de bananal
fumigado. O efeito sobre a média municipal é esse valor multiplicado pela
fração de nascimentos expostos. No Equador, 1,8\% dos nascimentos
geocodificados estão a até 150 metros de um bananal, em mais de 186 mil hectares
de plantações. Essa densidade, de cerca de 0,03 nascimento exposto por hectare
e por ano, transposta aos mil hectares da Chapada, dá cerca de trinta
nascimentos expostos por ano, de 1.068 em Limoeiro do Norte e
Quixeré, ou de 2.078 com Russas.\footnote{IBGE, Estatísticas do Registro Civil,
tabela SIDRA 2609, média de 2010 a 2018; os vinte e um municípios de controle
somam 8.874 nascimentos por ano.} A fração exposta fica entre 1,4\% e 2,8\%, e o
efeito esperado na média, no máximo entre um e dois gramas, pois o valor de
referência é o da exposição intensa.

Para o contraste entre a variação de tratados e de controles do pré-período ao
pós-período,
\[
EMD = (z_{1-\alpha/2} + z_{1-\beta})
\sqrt{\sigma^2\Big(\tfrac{1}{N_{T}^{\text{pré}}} + \tfrac{1}{N_{T}^{\text{pós}}} +
\tfrac{1}{N_{C}^{\text{pré}}} + \tfrac{1}{N_{C}^{\text{pós}}}\Big) + 2\,\tau^2\Big(\tfrac{1}{J_T}+\tfrac{1}{J_C}\Big)},
\]
em que $\sigma$ é o desvio-padrão individual do peso ao nascer, os $N$ são
números de nascimentos, $\tau$ é o desvio-padrão dos choques comuns aos
nascimentos de um município e $J$ é o número efetivo de municípios. A
Tabela~\ref{tab:emd} dá o valor com $\tau = 0$, que é um limite inferior, para
$\alpha = 5\%$, poder de 80\% e $\sigma = 500$ gramas. A coluna de base de
um ano corresponde ao $ATT(t)$ de um período; a do pré-período agrupado, a um
contraste agregado entre antes e depois.

\begin{table}[ht]
\centering\small
\caption{Limite inferior do efeito mínimo detectável sobre o peso ao nascer
(gramas)}
\label{tab:emd}
\begin{tabular}{@{}lrrr@{}}
\toprule
Grupo tratado & Nasc./ano & Pré 2010--2018 & Base de um ano \\
\midrule
Limoeiro do Norte e Quixeré & 1.068 & 23,9 & 49,0 \\
Com Russas                   & 2.078 & 18,0 & 36,9 \\
\bottomrule
\end{tabular}
\fonte{Elaboração própria. Considera só o erro amostral, com seis anos de
pós-período; $\sigma$ e $\tau$ são hipóteses de planejamento.}
\end{table}

O efeito esperado é de dez a trinta vezes menor que o limite inferior do efeito
mínimo detectável, e nenhuma das especificações tem poder para detectá-lo na
média municipal. A avaliação entrega, ainda assim, um limite para o efeito
médio, a medida de quanto avião a lei removeu, a fração da população ao alcance
da deriva e o comportamento dos demais elos da cadeia. O cálculo será refeito
com os microdados do pré-período. O intervalo de confiança decide: o efeito
mínimo detectável é instrumento de planejamento e não rebaixa nem promove,
depois, um resultado.

\section{Inferência}

Com poucos municípios tratados, o erro-padrão agrupado convencional não é
confiável. A inferência usa os métodos de \citeonline{conley2011} e de
\citeonline{ferman2019few} e testes de permutação entre os municípios de
controle. A sensibilidade a violações de tendências paralelas segue
\citeonline{rambachan2023credible}, e os desfechos secundários recebem correção
para comparações múltiplas.

\section{Segundo ensaio: custo-benefício e escolha do instrumento}
\label{sec:cb}

O segundo ensaio usa os resultados do primeiro para duas perguntas. A primeira é
de custo-benefício: o ganho em saúde, valorado pelo custo da doença
\cite{soares2009socialcost}, com preços reestimados para os dois estados, contra
a variação de área, produção e emprego e a diferença de custo entre as
tecnologias de aplicação. O resultado será apresentado com a pergunta invertida:
que efeito sobre o peso ao nascer seria necessário para a lei se pagar, e esse
efeito está dentro ou fora do intervalo estimado? A segunda é a escolha do
instrumento, entre proibir, taxar e delimitar zonas de exclusão, organizada por
\citeonline{weitzman1974}. Dois fatos brasileiros entram na discussão: a norma
federal já previa distâncias mínimas para a aplicação aérea, e o instrumento de
preço opera hoje como subsídio \cite{cunha2020incentivos}.
